\documentclass[10pt,a4paper]{amsart}
\usepackage[margin=19mm]{geometry}
\usepackage{amsmath,amssymb,mathtools}
\usepackage[scr=rsfs,cal=euler]{mathalfa}
\usepackage{xfrac}
\usepackage[unicode,hidelinks,bookmarks=false]{hyperref}
\newcommand{\ie}{i.e.\ }
\newcommand{\cf}{cf.\ }
\newcommand{\resp}{respectively\ }
\newcommand{\Subsection}{\subsection}
\providecommand{\nobreakdash}{\nobreak\hbox{-}}
\setlength{\emergencystretch}{2em}
\multlinegap0pt
\usepackage{tikz}
\usepackage{amssymb}
\usepackage{dsfont}
\usepackage{mathtools}
\usepackage{enumerate}
\usepackage{bm}
\usepackage[normalem]{ulem}
\newtheorem{lemma}{Lemma}
\newtheorem{theorem}{Theorem}
\newtheorem{proposition}{Proposition}
\def \cc {{\mathbb C}}
\DeclareMathOperator{\dom}{dom}
\def \rr {{\mathbb R}}
\DeclareMathOperator{\spec}{spec}
\newcommand{\specd}{\spec_\mathrm{disc}}
\newcommand{\spece}{\spec_\mathrm{ess}}
\title{The Landau--Dirac operator with shell interactions: \\self-adjointness and clustering}
\author{Badreddine Benhellal, Vincent Bruneau, Pablo Miranda}
\date{Source: arXiv:2608.20665v1 (CC BY 4.0)}
\begin{document}
\maketitle

\noindent\emph{Source and licence.} The Landau--Dirac operator with shell interactions: \\self-adjointness and clustering. Version arXiv:2608.20665v1 (CC BY 4.0), distributed by arXiv under the Creative Commons Attribution 4.0 International licence, http://creativecommons.org/licenses/by/4.0/, as recorded in the arXiv licence metadata for this exact version and shown on the abstract page https://arxiv.org/abs/2608.20665v1. Full licence notice: \texttt{licenses/arxiv-2608.20665-CC-BY-4.0.txt}.\par\smallskip

\noindent\textit{Editorial scope.} Counted unit: the article's proposition Spec\_H (source0.tex lines 1103--1108). The article's complete original proof of it is delivered with it (source0.tex lines 1109--1165, one \texttt{proof} environment of 57 lines). No mathematics is written, repaired or re-derived: the statement and the proof are the article's own lines, in the article's own notation, and no retained line is substituted or re-flowed.  Retained uncounted as the source-local context the proof needs: lines 658--662; lines 915--917; lines 923--925; lines 940--975; lines 1045--1050; lines 1059--1067.  Nothing was omitted from any retained span, so the package carries no omission record.  No retained line contains a citation, so the wrapper carries no bibliography and no bibliography entry.  No retained line contains a figure environment, an image-inclusion command or drawing code, so no image is delivered and the package declares no asset dependency.  Editorial formatting only, in addition: an AMS \texttt{amsart} preamble that restates the article's own declarations and macros used by the retained lines, namely the environments \texttt{lemma}, \texttt{theorem}, \texttt{proposition} on separate counters, together with the standard packages those lines need.  The retained environments print in this document as Lemma 1, Theorem 1, Proposition 1 in order of appearance, whereas the article prints the same items with the numbering of its own full document.  The complete native source of the article as distributed in the arXiv e-print for this exact version is bundled unchanged as \texttt{sources/208196/source0.tex}.
\begin{lemma}\label{Green} For any $u_\pm,v_\pm\in  H^1_{\sigma, A}(\Omega_\pm;\cc^2)$ there holds 
      \begin{align*}
    \langle -i\sigma\cdot\nabla_A u_\pm, v_\pm \rangle_{L^2(\Omega_\pm;\cc^2)}=\langle u_\pm, -i\sigma\cdot\nabla_A v_\pm \rangle_{L^2(\Omega_\pm;\cc^2)}+\langle \mp i\sigma\cdot\nu\gamma_D^\pm u_\pm, \gamma_D^\pm v_\pm \rangle_{L^2(\Sigma;\cc^2)}.
\end{align*}
\end{lemma}

     \begin{align}\label{assumption1}
     |\epsilon(s)|\neq |\tau(s)|\,\,\, \text{for any }\, s\in\Sigma.
     \end{align}

\begin{align}\label{TC}
    \frac{1}{2}(\epsilon I_2+\tau\sigma_3)(\gamma_D^+u_+ +\gamma_D^-u_-) +i\sigma\cdot\nu(\gamma_D^+u_+-\gamma_D^-u_-)=0,
 \end{align}

\begin{lemma}\label{TC_alter} Let $\epsilon,\tau\in C^1(\Sigma)$ satisfy \eqref{assumption1} and set 
\begin{equation*}\label{Projection}
    P_\pm=\frac{1}{2}\left(I_2 \mp \frac{i}{2}\sigma\cdot\nu(\epsilon I_2+\tau\sigma_3) \right).
\end{equation*}
Then, the following is true:  
\begin{itemize}
    \item[(i)] For those $s\in\Sigma$ such that   $\epsilon^2(s)-\tau^2(s)=-4$ we have that  $P_+$ and $P_-$ are projectors (i.e., $P_+ +P_-=I_2$ and $P_\pm^2= P_\pm$) and satisfy 
    \begin{align}\label{P_pm}
   P_\pm^\ast=\frac{1}{2}\left(I_2 \pm \frac{i}{2}\sigma\cdot\nu(\epsilon I_2-\tau\sigma_3) \right)\quad\text{and}\quad   P_\pm^\ast(i\sigma\cdot\nu)=(i\sigma\cdot\nu)P_\mp.
    \end{align}
  In particular, if  $\epsilon^2-\tau^2\equiv-4$ on $\Sigma$, then  for any $u=u_+\oplus u_-\in\dom D_{\epsilon,\tau}$ the transmission condition \eqref{TC} is equivalent to
    \[
    P_+ \gamma_D^+u_+=0 \quad\text{and}\quad  P_- \gamma_D^-u_-=0,
    \]
  and $D_{\epsilon,\tau}$ decomposes into the direct sum $D^{\Omega_+}_{\epsilon,\tau}\oplus D^{\Omega_-}_{\epsilon,\tau}$, where $D^{\Omega_\pm}_{\epsilon,\tau}$ are the Dirac operators defined by
    \begin{align}\label{H_Omega}
		\begin{split}
			D^{\Omega_\pm}_{\epsilon,\tau}u_\pm =& D u_\pm , \\
			\dom D^{\Omega_+}_{\epsilon,\tau} =& \{ v \in  H^1(\Omega_+,\cc^2): P_+ \gamma_D^+v=0  \},\\
            \dom D^{\Omega_-}_{\epsilon,\tau} =& \{ v \in H^1_{\sigma, A}(\Omega_-,\cc^2): P_- \gamma_D^-v=0  \}.
		\end{split}
	\end{align}
    \item[(ii)] For those $s\in\Sigma$ such that   $\epsilon^2(s)-\tau^2(s)\neq-4$ we have that   $P_+$ and $P_-$ are invertible with 
    \[
    (P_\pm)^{-1}=\frac{16}{\epsilon^2-\tau^2+4}P_{\mp}.
    \]
    In this case, for any $u=u_+\oplus u_-\in\dom D_{\epsilon,\tau}$ the transmission condition \eqref{TC} reads as follows
    \[
  Q_\pm \gamma_D^\mp u_\mp(s)=  \gamma_D^\pm u_\pm(s),
    \]
    where $Q_\pm$ are the invertible matrices
    \[
    Q_\pm= \frac{4}{\epsilon^2-\tau^2+4}\left(\frac{4 -\epsilon^2+\tau^2}{4}I_2 \pm i\sigma\cdot\nu(\epsilon I_2+\tau\sigma_3)) \right).
    \]
\end{itemize}
\end{lemma}

\begin{align}\label{Traceeigen}
\begin{split}
\frac{1}{2}(\epsilon I_2+\tau\sigma_3)(\gamma_D^+ u_+ +\gamma_D^- u_-)&=(\epsilon I_2+\tau\sigma_3)\mathscr{C}_zg,\\
 i\sigma\cdot\nu(\gamma_D^+ u_+ -\gamma_D^- u_-)&=g.
 \end{split}
\end{align}

\begin{theorem}\label{main_th1}Let $\epsilon,\tau\in C^1(\Sigma)$ satisfy $(\epsilon^2-\tau^2)(\Sigma) \cap\{0, 4\}=\emptyset$. Then for any $z \in \rho(D_0)\cap\rho(D_{\epsilon, \tau})$ the operator $ \Pi_{z}$ is boundedly invertible in $L^2(\Sigma; \mathbb{C}^2)$, and the resolvent formula
			\begin{equation}\label{Krein}
				(D_{\epsilon, \tau} - z)^{-1} = (D_0 - z)^{-1} - \Phi_{z} \Pi_{z}^{-1} \Phi_{\Bar{z}}^*,
\end{equation}				
holds. Moreover, $D_{\epsilon,\tau}$ is self-adjoint and there holds 
 \[
 \spece (D_{\epsilon,\tau})=\spece (D_0)=\spec (D_0).
 \]			
\end{theorem}

\begin{proposition}\label{Spec_H} Let $\epsilon,\tau\in C^1(\Sigma)$ satisfy $\epsilon^2(s)-\tau^2(s)= -4$ for all $s\in\Sigma$. Then, the spectrum of $D_{\epsilon,\tau}^{\Omega_+}$ is purely discrete and  $\spece(D_{\epsilon,\tau}^{\Omega_-})=\spece(D_0)$. In addition, if $\tau(s)\geq2$ 
holds for each $s\in\Sigma$, then we have
\[
\specd(D_{\epsilon,\tau}^{\Omega_+})\subset \rr\setminus [-m,m] \quad \text{and}\quad \specd(D_{\epsilon,\tau}^{\Omega_-})\subset \rr\setminus [-m,m).
\]
\end{proposition}

\begin{proof} Clearly $D_{\epsilon,\tau}^{\Omega_+}$ has a compact resolvent implying that its spectrum is purely discrete. Consequently, Theorem \ref{main_th1} yields that $\spece(D_{\epsilon,\tau}^{\Omega_-})=\spece(D_{\epsilon,\tau})=\spece(D_0)$.

We are now going to prove that $\specd(D_{\epsilon,\tau})\subset \rr\setminus [-m,m)$ holds for $\tau\geq2$ .  For any $u=u_+\oplus u_-\in\dom D_{\epsilon,\tau}$ we have 
\[
\|D_{\epsilon,\tau}u \|^2_{L^2(\rr^2;\cc^2)}= \|D_{\epsilon,\tau}^{\Omega_+}u_+ \|^2_{L^2(\Omega_+;\cc^2)}+\|D_{\epsilon,\tau}^{\Omega_-} u_- \|^2_{L^2(\Omega_-;\cc^2)},
\]
 and using Lemma \ref{Green}, we compute    
 \begin{align}\label{Quadratic1}
 \begin{split}
 \|D_{\epsilon,\tau}^{\Omega_\pm}u_\pm \|^2_{L^2(\Omega_\pm;\cc^2)}=& \langle D u_\pm, D u_\pm\rangle_{L^2(\Omega_\pm;\cc^2)}\\
 =&  \|\sigma\cdot\nabla_Au_\pm \|^2_{L^2(\Omega_\pm;\cc^2)} + m^2  \| u_\pm \|^2_{L^2(\Omega_\pm;\cc^2)}\\
 &+m\langle \mp i\sigma\cdot\nu \gamma_D^\pm u_\pm, \sigma_3 \gamma_D^\pm u_\pm\rangle_{L^2(\Sigma;\cc^2)} .  
 \end{split}
 \end{align}
By Lemma \ref{TC_alter}-(i), we have 
 \begin{align}\label{Quadratic2}
 \gamma_D^\pm u_\pm= \pm\frac{i}{2}\sigma\cdot\nu(\epsilon I_2+\tau\sigma_3)\gamma_D^\pm u_\pm.
 \end{align}
 Thus, if write  $u_\pm =(v_\pm,w_\pm)^\top$ then  \eqref{Quadratic2} is equivalent to
 \[
 \mp i\sigma\cdot\nu \gamma_D^\pm u_\pm= \begin{pmatrix}
     \frac{\epsilon+\tau}{2} & 0\\
     0 & \frac{\epsilon-\tau}{2}
 \end{pmatrix} \gamma_D^\pm u_\pm  = \frac{1}{2} \begin{pmatrix}(\epsilon+\tau)\gamma_D^\pm v_\pm\\
(\epsilon-\tau) \gamma_D^\pm w_\pm
 \end{pmatrix}.
 \]
 This implies 
 \begin{align*}
 \langle \mp i\sigma\cdot\nu \gamma_D^\pm u_\pm, \sigma_3 \gamma_D^\pm u_\pm\rangle_{L^2(\Sigma;\cc^2)}= \frac{1}{2}\int_\Sigma \left[  (\epsilon+\tau) |v_\pm |^2 + (\tau-\epsilon) |w_\pm |^2 \right]\mathrm{d}s.
 \end{align*}
  Plugging this into  \eqref{Quadratic1}  yields 
 \begin{align}\label{Quadratic3}
 \begin{split}
 \|D_{\epsilon,\tau}u \|^2_{L^2(\rr^2;\cc^2)} =&  \|\sigma\cdot\nabla_A u_+ \|^2_{L^2(\Omega_+;\cc^2)} + \|\sigma\cdot\nabla_A u_- \|^2_{L^2(\Omega_-;\cc^2)} + m^2 \|u \|^2_{L^2(\rr^2;\cc^2)} \\
 &+ \frac{m}{2} \int_\Sigma (\epsilon+\tau)(|v_+ |^2 + |v_- |^2) \mathrm{d} s   + \frac{m}{2} \int_\Sigma (\tau-\epsilon)( |w_+ |^2 + |w_- |^2)\mathrm{d}s.
 \end{split}
 \end{align}
Since $\tau \geq2$ it follows that $\tau>|\epsilon|$, and  \eqref{Quadratic3} entails that
 \begin{align*}
 \|D_{\epsilon,\tau}u \|^2_{L^2(\rr^2;\cc^2)} \geq m^2 \|u \|^2_{L^2(\rr^2;\cc^2)}
 \end{align*}
 showing that $\spec(D_{\epsilon,\tau})\subset \rr\setminus (-m,m)$. In view of \eqref{Quadratic3}, if $-m$ is an eigenvalue of $D_{\epsilon,\tau}$, then its associated eigenfunction $u$ satisfies $\gamma_D^\pm u_\pm=0$. However, $u=\Phi_{-m} g$ with $g=i\sigma\cdot\nu(\gamma_D^+ u_+-\gamma_D^- u_-)$ by \eqref{Traceeigen}, and we conclude that $u$ vanishes identically on $\rr^2$. Therefore, $-m\notin\spec(D_{\epsilon,\tau})$ which leads to $\specd(D_{\epsilon,\tau})\subset \rr\setminus [-m,m)$.  

Similarly, if  $m\in \specd (D_{\epsilon,\tau}^{\Omega_+})$ then by  \eqref{Quadratic3} its associated eigenfunction $u_+\in \dom D_{\epsilon,\tau}^{\Omega_+}\setminus\{0\}$ satisfies 
\begin{align*}
\gamma_D^+ u_+=0 \,\,\text{on }\,\Sigma \quad\text{and}\quad \sigma\cdot\nabla_A u_+=0 \,\,\text{in }\,\Omega_+.
\end{align*}
Hence, $u_+ =(v_+,w_+)^\top\in H^1_0(\Omega_+;\cc^2)$ and applying $\sigma\cdot\nabla_A$ to the second equality above gives
\begin{align*}
0=(\sigma\cdot\nabla_A)^2 u_+\;\;
&=\; \begin{pmatrix} -\nabla_A^2 - b & 0 \\ 0 & -\nabla_A^2 + b \end{pmatrix}\begin{pmatrix} v_+\\
  w_+
 \end{pmatrix}.
\end{align*}
This contradicts the fact that the infimum of the spectrum of the magnetic Schr\"odinger operator with Dirichlet boundary condition is larger than the lowest Landau level $b$. Therefore,  $m\notin \specd (D_{\epsilon,\tau}^{\Omega_+})$ and the proposition is proved.
\end{proof}

\end{document}
